\section{Introduction}\label{sec:intro}

Let $X$ be a smooth complex projective variety. The Hodge conjecture asserts
that the cycle class map
\[
  \cl^{p}_{X}\colon \CH^{p}(X)_{\QQ}\longrightarrow \Hdg^{p}(X)
  \coloneqq H^{2p}(X,\QQ)\cap H^{p,p}(X)
\]
is surjective for every $p$. Beyond degree two the best understood Hodge
classes are those of abelian varieties, and the first of them that divisors
do not produce are the Weil classes, which Weil exhibited on abelian
varieties with an action of an imaginary quadratic field whose two
eigenspaces on the tangent space have the same dimension \cite{Weil77}; the
construction works for every CM field $F$. Markman proved that the Weil
classes are algebraic on the abelian sixfolds of split Weil type for every
imaginary quadratic field, and, through a product argument of Schoen
\cite[Proposition~10]{Sch98}, on every abelian fourfold of Weil type
\cite{Mar25a,Mar25b}. His sheaves are built from secant sheaves on an abelian
variety of half the dimension, and they extend along the family because they
satisfy an equivariant form of semiregularity.

The companion paper \cite{paperI} studies the set of members of a family on
which a flat Hodge class is algebraic, its \emph{algebraic locus}, and
determines it for the Weil classes of every CM field and every dimension at
once: for a CM field $F$, an integer $n\ge2$ and a discriminant $\delta$, the
abelian varieties of $(F,n,\delta)$-Weil type form one family over a
connected period domain $\cD_{F}$, and the locus $\Sigma_{F}$ of members on
which the Weil classes are algebraic contains an explicit point of complex
multiplication, is dense, is a countable union of closed algebraic strata of
the Shimura variety of the family, and is either all of $\cD_{F}$ or meagre;
so the Weil classes are algebraic on every member exactly when $\Sigma_{F}$
is closed (\cite[Theorem~1.1]{paperI}, which we call the main theorem throughout). The
companion paper \cite{paperII} determines what closing $\Sigma_{F}$ requires
of a single object at a single point: a perfect complex at a base point
whose Chern character is a nonzero multiple of a Weil class plus a
polynomial in the polarisation and whose semiregularity map is injective
makes $\Sigma_{F}$ everything (\cite[Theorem~2.3]{paperII}), the criterion is
never met when the Chern character is exactly a Weil class
(\cite[Theorem~3.5]{paperII}, \cite[Theorem~5.4]{paperII}), and with powers of the
polarisation added it becomes one number (\cite[Theorem~3.9]{paperII}). The
notation of those papers is kept here.

This paper places the Weil classes inside the whole conjecture. It
determines the reach of the obstructions of \cite{paperI}, shows that the
Hodge conjecture is two statements about algebraic loci, proves each on the
families where it can be proved, proves the conjecture itself on a family of
varieties of every dimension, and studies the first instance of the
statement on abelian varieties, the exceptional classes of a Mumford
fourfold.

\subsection{Diagonal complete intersections}

The paper also proves the Hodge conjecture itself on a family of varieties
of every dimension. Let $X\subset\PP^{N}$ be a complete intersection of $c$
diagonal hypersurfaces $\sum_{i}a_{ki}x_{i}^{d}=0$ of degree $d$. It is
\emph{of Vandermonde type} if $a_{ki}=\mu_{i}\sum_{l=1}^{c}g_{kl}\lambda_{i}^{l-1}$
for distinct numbers $\lambda_{0},\dots,\lambda_{N}$, nonzero numbers
$\mu_{i}$ and an invertible matrix $(g_{kl})$; for $c\ge2$ this says that the
columns of $(a_{ki})$ are distinct points of a rational normal curve, and
every smooth complete intersection of two diagonal hypersurfaces of the same
degree is of this type. Such an $X$ is isomorphic to a variety
$X_{d,N-c}(\lambda)$ that depends only on $\lambda$ (\Cref{thm:vandermonde}),
and \emph{very general} refers to $\lambda$. For $d=2$ we call $X$ a
\emph{Vandermonde quadric intersection}.

\begin{theorem}[Diagonal complete intersections]\label{thm:introvandermonde}
Every rational Hodge class is algebraic on the very general Vandermonde
quadric intersection, in every dimension and in every projective space. The
same holds for the very general complete intersection of Vandermonde type of
diagonal cubics or quartics of dimension at most seven, or of diagonal
sextics of dimension at most five, and for every smooth complete
intersection of two diagonal cubics, of two diagonal quartics or of two
diagonal sextics in $\PP^{6}$.
\end{theorem}

The first two statements are \Cref{thm:vgvandermonde}(iii) and the third is
\Cref{cor:twodiagonal}(ii). The variety $X$ is a quotient of a power of a
generalised Fermat curve, so its Hodge classes come from abelian subvarieties
of the Jacobian of that curve. For quadrics the monodromy of the curve along
a chain of vanishing cycles is the full symplectic group on every eigenspace,
so every Hodge class comes from powers of polarisations and no theorem on
abelian varieties of higher dimension is needed. For cubics, quartics and
sextics we compute the monodromy by an induction on the number of branch
points (\Cref{prop:cyclicmonodromy}), and the classes it leaves are divisor
classes, Weil classes of abelian fourfolds of Weil type, and Weil classes of
abelian sixfolds of split Weil type, all algebraic by the theorems of
Lefschetz and Markman \cite{Mar25a,Mar25b}. In odd dimension every Hodge
class is a multiple of a power of the hyperplane class. So the bound seven
for cubics and quartics is the point where abelian varieties of dimension
eight appear, and the bound five for sextics comes from abelian sixfolds of
non-split Weil type, which occur in dimension six.



\subsection{Weil classes and the Hodge conjecture}

\Cref{sec:scope,sec:beyond,sec:families} place the Weil classes inside the
conjecture. Two
statements carry all of it. The first, (F2), asks that every Hodge class on an
abelian variety outside the subring generated by divisor classes and Weil
classes of quotients be algebraic; multiplying by an exceptional class on the
square of a Mumford fourfold moves every Hodge class of every abelian variety
outside that subring, so (F2) is equivalent to the Hodge conjecture for all
abelian varieties (\Cref{prop:f2isab}). The second, (F3$'$), asks that every
Hodge class on every smooth projective variety be algebraic modulo the images
of Hodge classes of abelian varieties under algebraic correspondences
(\Cref{def:f3prime}). The Hodge conjecture is equivalent to (F2) and (F3$'$)
together (\Cref{prop:f3prime}, \Cref{thm:final}).

We prove (F3$'$) for every variety whose cohomology is reached from abelian
varieties by correspondences, a class closed under products, surjective
images, blow-ups, projective bundles and Hilbert schemes of points on
surfaces, which contains every Fermat hypersurface and every K3 surface whose
Kuga--Satake correspondence is algebraic (\Cref{prop:f3prime},
\Cref{prop:aclosure}); for every variety of dimension at most three, and on
a variety of dimension $k$ in the degrees $0$, $2$, $2k-2$ and $2k$; for
every fourfold and fivefold whose zero-cycles are supported in dimension at
most three, among them every uniruled fourfold and every rationally connected
fivefold (\Cref{prop:f3primesmall}, \Cref{prop:f3primefibration}); for every
fourfold and fivefold dominated by a variety of that class
(\Cref{prop:f3primedominant}); for every smooth ample divisor of odd
dimension in a variety of that class, and in even dimension outside the
middle degree and for a very general member (\Cref{prop:f3primeample});
and, by a toric compactification whose boundary is again of the same kind, for
every hypersurface of simplex type, so for every smooth Delsarte hypersurface
of every dimension and every cyclic cover branched along one, whose Hodge
conjecture then follows from that of Fermat varieties of one degree
(\Cref{thm:simplextype}, \Cref{cor:delsarteall}). A complete intersection
of diagonal hypersurfaces whose coefficients are of Vandermonde type, which
holds for every smooth complete intersection of two of them, is a quotient
of a power of a generalised Fermat curve, so it satisfies (F3$'$), and its
Hodge conjecture reduces to abelian varieties attached to the characters of
the curve; for two diagonal hypersurfaces of degree $3$, $4$ or $6$ in
$\PP^{6}$ these have dimension at most five, and the Hodge conjecture holds,
the exceptional classes being Weil classes of abelian fourfolds of Weil type
(\Cref{thm:vandermonde}, \Cref{cor:twodiagonal}). At a very general
configuration the monodromy of the curve leaves only the Hodge classes that
come from polarisations and from Weil classes; we compute that monodromy for
cyclic covers of degree three, four and six by an induction on the number of
branch points (\Cref{prop:cyclicmonodromy}), so the Hodge conjecture holds
for the very general such intersection of quadrics in every dimension, of
cubics and quartics up to dimension seven, and of sextics up to dimension
five, in every projective space (\Cref{thm:vgvandermonde}). On the powers, Hilbert
schemes and moduli spaces of a K3 surface whose transcendental field is a CM
field, and on the fourfolds of K3$^{[2]}$ type with field $\QQ$ or a CM
field, the Hodge conjecture itself holds (\Cref{prop:k3powers},
\Cref{prop:k3ntype}).

The first instance of (F2) is a pair of classes on the square of a Mumford
fourfold, and \Cref{sec:beyond} studies it in detail. The two classes are
the real multiplication by a totally real cubic field on the primitive part of
$H^{2}$, and they are algebraic once the Kuga--Satake class of one K3 surface
of Picard number thirteen is algebraic (\Cref{prop:mumfordrm},
\Cref{thm:mumfordks}); they are algebraic at every CM point
(\Cref{thm:mumfordcm}); they are rigid, their Hodge locus being the compact
Mumford curve (\Cref{thm:mumfordrigid}); and they are represented by
algebraic cycles at every closed point of every reduction of the curve modulo
a prime, so that their algebraicity on every member is a bound on the Hilbert
data of those cycles (\Cref{prop:modp}, \Cref{thm:modpdensity}). For an
abelian scheme over a curve whose invariant classes are algebraic the
Lefschetz standard conjecture holds on the total space
(\Cref{thm:lefschetzfamily}), and over a compact Shimura curve the total space
of a Mumford family satisfies the Hodge conjecture
(\Cref{prop:mumfordfivefold}).

The structure that \cite[Theorem~1.1]{paperI} finds in one family is present in every
family that carries (F2) or (F3$'$) (\Cref{sec:families}). Every Hodge class
of every abelian variety lies in an abelian scheme with a member of CM type,
and (F2) is equivalent to the statement that in every such scheme the
algebraic locus of a flat Hodge class that is algebraic at a CM member has an
interior point (\Cref{prop:cmfamily}, \Cref{thm:f2propagation}); the base
points of the Weil families are the seeds, through Andr\'e's theorem on the
Hodge classes of abelian varieties of CM type. (F3$'$) is equivalent to the
absence of Hodge classes off the part of the cohomology reached from abelian
varieties by algebraic correspondences (\Cref{prop:abpart}). On the squares of
K3 surfaces with real multiplication, where (F3$'$) is first open, the
algebraic locus of the real multiplication contains every CM point, is dense,
is stable under rational Hodge isometries, and is either the whole family or
meagre (\Cref{thm:rmlocus}). In each case the property still missing is the
interior point (\Cref{thm:twostand}).


Finally, the implications of the three papers and of the literature form a
finite rule set, and its consequences can be determined by hand. They give
the minimal sets of statements from which the conjecture follows
(\Cref{thm:closure}, \Cref{rem:f2rule}): (F3$'$) together with the
Lefschetz standard conjecture, with the variational Hodge conjecture for
algebraic classes, or with (F2), besides two sets that are each equivalent to
the conjecture.%
\footnote{The Lefschetz standard conjecture asks that the inverse of the
hard Lefschetz isomorphism $L^{\dim X-k}\colon H^{k}(X,\QQ)\to
H^{2\dim X-k}(X,\QQ)$ be induced by an algebraic correspondence. The
variational Hodge conjecture asks that a flat section of Hodge classes over
the base of a smooth projective family that is algebraic at one fibre be
algebraic at every fibre.}
\Cref{sec:final} collects the results in one theorem.

Every proof is written to be checked by hand. Where an argument rests on a
finite calculation, the calculation is given where it is used, as an explicit
identity, a determinant or a case analysis. The archive \cite{BMB26} holds
programs that repeat much of this arithmetic and a Lean~4 certificate of
part of it; no statement of the paper rests on them.
